\section{A nearly wire-preserving 3-SAT encoding}\label{sec:normalization}

The normalized contractor parameter is a wire count. An application to
SAT must pay for its actual Boolean network; replacing that parameter by
the number of variables without a reduction is invalid. This section
gives a direct, exact normalization with a precise wire formula.

\subsection{Clause and equality tensors}

First delete tautological clauses, repeated literals within a clause,
and duplicate clauses. An empty clause decides unsatisfiability. A formula
with no remaining clauses has $2^n$ satisfying assignments. Rename the
remaining variables by their active indices and record the number $n_u$
of unused declared variables. These operations preserve the count on
active variables, with an eventual factor $2^{n_u}$.

For each remaining clause, use its Boolean truth table on its one, two,
or three occurrence wires. For a variable with $d$ occurrences, an
equality tensor on the variable side is one exactly when all its ports
carry the same bit. The direct incidence network has $L$ wires, where
$L$ is the normalized literal count. It is already exact, but a variable
of large occurrence may violate the root degree premise.

\begin{theorem}[Bounded-degree equality normalization]\label{thm:equality}
For any integer $B\ge3$, a normalized 3-CNF with active occurrence
counts $d_v\ge1$ has a bipartite Boolean network of maximum degree
at most $B$ and exactly
\begin{equation}\label{eq:wires}
 M_B=L+2\sum_v\floor{\frac{(d_v-3)_+}{B-2}}
 \le L\left(1+\frac2{B-2}\right)
\end{equation}
wires. Its contraction is the exact number of satisfying assignments
on the active variables. Its description and entry queries are
polynomial-time computable. Unused variables contribute the separate
factor $2^{n_u}$.
\end{theorem}

\begin{proof}
Replace the equality star of variable $v$ by a path of
\[
 k_v=\max\left\{1,\ceil{\frac{d_v-2}{B-2}}\right\}
\]
equality nodes, all on the variable side. A path with $k_v$ such nodes
has total external-port capacity
\[
 k_vB-2(k_v-1)=k_v(B-2)+2\ge d_v.
\]
Assign its $d_v$ occurrence ports along the path without exceeding these
capacities. Subdivide every path edge by one binary equality node on the
clause side. Thus the graph stays bipartite. Each subdivision contributes
two internal wires, so this variable adds $2(k_v-1)$ wires beyond its
original $d_v$ clause incidences. The identity
\[
 k_v-1=\floor{\frac{(d_v-3)_+}{B-2}}
\]
gives the exact formula. Summing the elementary upper bound
$k_v-1\le d_v/(B-2)$ proves its inequality.

Because the equality nodes form a connected tree, an assignment of
external occurrence bits extends to its internal wires if and only if
all external bits agree. If it extends, the extension is unique. There
is therefore neither an unwanted multiplicity nor an extra restriction.
Summing over occurrence wires counts exactly the satisfying assignments
of the active formula. Every node has degree at most $B$, and each
entry query tests either a clause of size at most three or equality of
at most $B$ supplied bits. Construction and entry evaluation are
polynomial in the network description.
\end{proof}

The tree size is optimal within this particular family of degree-$B$
equality trees with binary subdivisions: any connected tree of $k$
variable-side equality nodes has only $k(B-2)+2$ available external
ports. This is a scoped combinatorial observation, not a lower bound
against arbitrary SAT encodings or tensor factorizations.

\subsection{Application to the low-degree contractor}

Choose
\[
 B=\max\{3,\ceil{(\log_2(L+2))^2}\},\qquad s=M_B+1.
\]
Then $s=(1+o(1))L$, and $B\le s^{1/12}$ for all sufficiently large
$L$. Small cases may be evaluated directly. Set the initial precision
to $P_0=n_a+1$, where $n_a\le L$ is the number of active variables,
and set $\ell=1$. The true count is at most $2^{n_a}$, so its residue
modulo $2^{n_a+1}$ uniquely recovers it. In particular $P_0\le s$,
as required by the interface. Multiply the recovered count by $2^{n_u}$
to restore unused variables. Writing this binary output costs at most
a polynomial in the original $n+m$.

All original tensor entries are computable in polynomial time, which is
stronger than the premise of Theorem~\ref{thm:contractors}. Since
\[
 \frac{s}{\sqrt{\log_2s}}=(1+o(1))\frac L{\sqrt{\log_2L}},
\]
the low-degree theorem gives Theorem~\ref{thm:sat}. Its decision answer
is simply whether the exact count is nonzero. With $L\le3m$, monotonicity
of $x/\sqrt{\log_2x}$ for sufficiently large $x$ gives the displayed
clause and bounded-density consequences.

\subsection{Comparison and limits of the normalization}

A trivalent equality tree has up to $3L$ wires under the same bipartite
subdivision convention. Allowing the degree cap to grow slowly makes
the added internal wires $o(L)$ instead of a constant multiple of $L$.
The construction is simple, but it is useful because the low-degree
contractor permits a much larger eventual degree than a constant.

This factor saving should not be presented as an additional factor-three
improvement over the strongest composition of all earlier ideas.
The direct-incidence/high-occurrence conditioning method already appears
in \cite{balanced}. In that method, condition on every variable whose
original occurrence exceeds $B$. There are at most $L/B$ such variables,
so it uses at most $2^{L/B}$ branches. Each residual incidence network
has at most $L$ wires and degree at most $B$; simplify satisfied clauses
and restore any now-unused variables when counting. With the same
$B=\Theta(\log^2L)$, the branching exponent is
$o(L/\sqrt{\log L})$. Its low-degree schedule can use the original
budget $s=L+1$; a smaller residual need not be treated as a fresh
instance with an unjustified relative degree bound.

Thus the earlier low-degree contractor already implies a fair leading
cached benchmark of $13.3796490663\ldots$ per literal. Our new
$10.3367239856\ldots$ coefficient beats that benchmark. The equality
encoding offers a single count-preserving network and an exact wire
formula, avoiding the outer branching construction.

Only original clause and equality tensors have been split. The proof
does not factor arbitrary derived tensors created by the compressed
dense-formula reduction. Such tensors may lack an equality-tree
representation. Consequently this normalization alone does not remove
the $\log\log n$ factor in the unrestricted variable-based bound.

The supplied compiler implements this construction, validates its degree,
port and bipartiteness invariants, and provides brute-force semantics for
small instances. The companion verifier compares exact tensor sums with
direct formula counts and checks the unique-extension property separately.
It is a normalization implementation, not the complete asymptotic solver.
